% Generated draft: requires mathematical and source review.
\documentclass[11pt]{article}
\usepackage[margin=25mm]{geometry}
\usepackage{fontspec}
\setmainfont{Latin Modern Roman}
\usepackage{amsmath,amssymb,mathtools,unicode-math}
\setmathfont{Latin Modern Math}
\usepackage{enumitem,needspace,xurl,hyperref}
\setlength{\parindent}{0pt}
\setlength{\parskip}{5pt}
\setlength{\emergencystretch}{3em}
\newcommand{\N}{\mathbb{N}}
\newcommand{\Z}{\mathbb{Z}}
\newcommand{\Q}{\mathbb{Q}}
\newcommand{\R}{\mathbb{R}}
\newcommand{\C}{\mathbb{C}}
\newcommand{\F}{\mathbb{F}}
\newcommand{\A}{\mathbb{A}}
\newcommand{\PP}{\mathbb{P}}
\newcommand{\E}{\mathbb{E}}
\newcommand{\eps}{\varepsilon}
\newcommand{\dd}{\,\mathrm{d}}
\begin{document}
\section*{UnsolvedMath ID: 30002565}
\textbf{Source problem code:} catalogue entry
\subsection{Definitions and mathematical statement}
Let \(\Z^2=\{(m,n):m,n\in\Z\}\) be the integer lattice with componentwise addition. The supplied record asks the following methodological question:
\begin{quote}
Can the finite-dimensional completely-positive-functional and separation argument used for free groups be extended to \(\Z^2\), so as to better understand this group?
\end{quote}
This question does not specify a mathematical conclusion to be proved. In particular, the supplied record does not identify the free-group argument, its hypotheses, the algebra and matrix-positive cones on which the completely positive functionals are defined, what is required to be finite-dimensional, or the objects and topology involved in separation. It also does not specify the ranks of the free groups under consideration or the intended conclusion for \(\Z^2\). Consequently, no fully quantified theorem or equivalence can be extracted without additional source information. The question above retains the methodological target; it does not assert any particular positivity or sums-of-squares characterization.
\subsection{Short English statement}
The question asks whether a method involving finite-dimensional completely positive functionals and separation, used for free groups, can also clarify the case of \(\Z^2\). The argument and desired conclusion are unspecified, so a precise theorem remains to be identified.
\subsection{Sources}
\begin{itemize}
\item Local UnsolvedMath registry: supplied record with permanent ID \(30002565\), title ``Completely Positive Separation for the Integer Lattice,'' source problem code \texttt{catalogue entry}, in the Oberwolfach Reports Open Problems collection. The recorded proposed year is 2014; no proposer is named.
\item Supplied report citation: \emph{Real Algebraic Geometry With A View Toward Systems Control and Free Positivity} (2015), Oberwolfach Reports, pp.~977--1045, DOI \texttt{10.4171/owr/2014/17}. Source URL and supplied literature-source link, labelled ``Oberwolfach Report 2014/17'': \url{https://doi.org/10.4171/owr/2014/17}. No passage, contribution, or theorem locating the referenced argument is supplied.
\item Literature assessment dated exactly \texttt{2026-08-21}: the supplied curation labels the problem open and reports that Scheiderer proves the relevant sums-of-squares characterization for \(\Z^2\), while the extracted methodological request is not a precise theorem and no definitive resolution of that request is identifiable. This is supplied dated curation, not an independently verified assessment. The record expressly requests independent verification of time-sensitive status and cited claims; it supplies no separate citation for Scheiderer's result.
\item The registry labels the statement \texttt{corrected\_verified} and reports manual semantic, typographic, and LaTeX review of a meaning-preserving cleanup. These are supplied editorial assertions and do not resolve the missing mathematical specifications.
\end{itemize}
\end{document}
